\documentclass[conference]{IEEEtran}
\usepackage[utf8]{inputenc}
\usepackage{graphicx}
\usepackage{amsmath}
\usepackage{amssymb}
\usepackage{listings}
\usepackage{xcolor}
\usepackage{booktabs}
\usepackage{hyperref}
\usepackage{array}
\usepackage{float}
\usepackage{caption}

% ---------- Listing caption fix ----------
\captionsetup[lstlisting]{
  position=top,
  skip=8pt,
  font=small,
  labelfont=bf
}
\renewcommand{\lstlistingname}{Listing}

\lstset{
  basicstyle=\ttfamily\footnotesize,
  breaklines=true,
  frame=single,
  backgroundcolor=\color{gray!10},
  keywordstyle=\color{blue},
  commentstyle=\color{green!50!black},
  stringstyle=\color{red!60!black},
  aboveskip=1.5em,
  belowskip=1.5em,
  captionpos=t,
  xleftmargin=0.5em,
  xrightmargin=0.5em,
  framexleftmargin=0.5em,
  framexrightmargin=0.5em,
}

\title{Tetris OS Simulator: A Custom Operating System\\Module Demonstration Using Game Engine Architecture}
\author{
  \IEEEauthorblockN{Mehak Jain, Anuj Kumar Singh}
  \IEEEauthorblockA{Operating System}
}

\begin{document}

\maketitle

% ============================================================================
\begin{abstract}
This paper presents the design, implementation, and evaluation of \textbf{Tetris OS Simulator} --- a fully-functional Tetris game engine built from scratch in C without standard library dependencies for core logic. The project simulates seven fundamental operating system modules: Process Management, Memory Management, File Systems, I/O Management, Error Handling \& Security, Networking, and User Interface. The engine runs in two modes --- a terminal-based interface using ANSI escape codes and a WebSocket-backed browser client using React. Six custom libraries (\texttt{memory.c}, \texttt{string.c}, \texttt{math.c}, \texttt{screen.c}, \texttt{keyboard.c}, \texttt{sound.c}) replace all standard library functionality, while a hardware abstraction layer confines POSIX networking to dedicated I/O modules. The system demonstrates arena-based dynamic memory allocation, persistent file-based leaderboards, real-time WebSocket communication with automatic reconnection, and comprehensive error handling across all subsystems.
\end{abstract}

\begin{IEEEkeywords}
Operating Systems, Custom Libraries, Memory Management, WebSocket, Tetris, Process Scheduling, Hardware Abstraction
\end{IEEEkeywords}

% ============================================================================
\section{Introduction}

\subsection{Problem Statement}
Standard game engines rely heavily on standard library functions (\texttt{malloc}, \texttt{strlen}, \texttt{sin}, etc.) and opaque frameworks. This project challenges that paradigm by implementing a complete game system where all core logic is built from first principles --- custom arithmetic, custom string operations, custom memory allocation, and custom I/O drivers. The complete source code and full commit history are publicly available at \url{https://github.com/Mehak261124/Tetris}.

\subsection{Objectives}
\begin{enumerate}
  \item Implement six custom foundational libraries replacing \texttt{<string.h>}, \texttt{<math.h>}, and \texttt{malloc/free}.
  \item Build a complete Tetris game engine demonstrating all seven OS module categories.
  \item Provide dual I/O modes: terminal (ANSI) and browser (WebSocket).
  \item Ensure robust error handling, edge-case coverage, and graceful degradation.
  \item Maintain persistent state across sessions (file system module).
  \item Support real-time client-server communication (networking module).
\end{enumerate}

\subsection{Scope}
The project covers the complete lifecycle of an OS-like system: boot (memory initialization), process scheduling (game loop), I/O device management (keyboard/screen/speaker), persistent storage (leaderboard files), and graceful shutdown (memory cleanup, socket release).

% ============================================================================
\section{System Architecture}

The system follows a layered architecture where each custom library represents an OS module:

\begin{lstlisting}[caption={System Architecture Layers}]
Application Layer  : ws_tetris.c / main.c
Core Library Layer : string.c | math.c | memory.c
HAL Layer          : screen.c | keyboard.c | sound.c
Host OS / Hardware : POSIX sockets | Terminal | Audio
\end{lstlisting}

\textbf{Design Principle:} Game logic never includes networking headers, POSIX headers, or standard library functions beyond \texttt{<stdio.h>} and \texttt{<stdlib.h>}. All I/O flows through the Hardware Abstraction Layer (HAL).

% ============================================================================
\section{Methodology: Seven OS Modules}

% ---- Module 1 ----
\subsection{Module 1: Process Management}
Process management determines \textit{who gets to run, when, and where}. The game loop acts as a cooperative scheduler executing three processes each tick:

\begin{table}[H]
\centering
\caption{Process Schedule}
\begin{tabular}{lll}
\toprule
\textbf{Process} & \textbf{Priority} & \textbf{Frequency} \\
\midrule
Input Polling    & Highest & Every tick (50ms) \\
Physics/Gravity  & Medium  & Level-dependent   \\
Rendering        & Lowest  & On state change   \\
\bottomrule
\end{tabular}
\end{table}

\textbf{Speed Curve (NES-style):}
\begin{itemize}
  \item Level 1--8: \texttt{drop\_ticks = max(13 - level, 5)} --- 600ms to 250ms
  \item Level 9--13: \texttt{drop\_ticks = max(8 - (level-8)/2, 2)} --- 150ms to 100ms
  \item Level 14+: \texttt{drop\_ticks = 1} --- 50ms (maximum speed)
\end{itemize}

\textbf{Lock Delay:} A 500ms timer prevents instant locking when a piece touches the ground. The timer resets on successful moves or rotations, capped at 15 resets per piece to prevent infinite stalling.

\textbf{DAS/ARR:} The web client implements Delayed Auto Shift (170ms delay) followed by Auto Repeat Rate (50ms interval) for smooth, responsive key repeat matching modern Tetris standards.

\textbf{Signal Handling:} \texttt{SIGINT} and \texttt{SIGTERM} are intercepted by a lightweight handler that sets a \texttt{volatile sig\_atomic\_t} flag. The game loop checks this flag once per frame and exits through the normal teardown path, ensuring \texttt{keyboard\_restore()} and \texttt{memory\_cleanup()} always run --- the terminal is never left in raw mode after a forced quit. \texttt{SIGWINCH} triggers a debounced resize handler: an 80\,ms \texttt{usleep()} wait discards all intermediate signals fired during continuous window drag, then a single \texttt{screen\_clear()} with the \texttt{\textbackslash{}033[3J\textbackslash{}033[2J\textbackslash{}033[H} sequence erases iTerm2's saved scrollback lines before redrawing the full frame cleanly.

% ---- Module 2 ----
\subsection{Module 2: Memory Management}
Memory management determines \textit{who owns which memory and how much}.

\textbf{Design:} First-Fit Free-List Allocator with Block Splitting \& Coalescing.

\begin{lstlisting}[language=C,caption={BlockHeader Structure}]
typedef struct BlockHeader {
    int                  size;
    int                  is_free;
    struct BlockHeader  *next;
} BlockHeader;
\end{lstlisting}

\textbf{Operations:}
\begin{itemize}
  \item \texttt{memory\_init(1\,048\,576)} --- acquires one 1\,MB slab via the single permitted \texttt{malloc()} call.
  \item \texttt{t\_alloc(size)} --- first-fit search with 8-byte alignment; splits oversized blocks to minimise waste.
  \item \texttt{t\_dealloc(ptr)} --- validates pointer ownership by walking the block list, marks free, then coalesces adjacent free blocks in a single pass.
  \item \texttt{memory\_cleanup()} --- releases the entire slab back to the host OS via the single permitted \texttt{free()} call.
\end{itemize}

\textbf{Safety guarantees:} 8-byte alignment prevents misaligned struct access; pointer validation rejects bad inputs without crashing; double-free is silently ignored; \texttt{NULL} is returned on allocation failure so callers can check before use.

% ---- Module 3 ----
\subsection{Module 3: File System}
The file system provides \textit{persistent storage} for leaderboard history across sessions.

\begin{itemize}
  \item \textbf{Shared leaderboard file} (\texttt{leaderboard.txt}) --- both the terminal build (\texttt{main.c}) and the WebSocket backend (\texttt{ws\_tetris.c}) read and write the same file using the same format, so scores are shared across both play modes.
  \item \textbf{Unified format:} \texttt{score|level|lines|name} (pipe-delimited, one entry per line, sorted descending by score). Top 5 entries stored in terminal mode, top 10 in web mode.
  \item \textbf{Parsing:} \texttt{t\_split()} from \texttt{string.c} tokenises each line --- no \texttt{strtok} or \texttt{sscanf}.
  \item \textbf{Edge cases:} Missing file starts with an empty table (not an error); malformed lines are silently skipped; write failures degrade gracefully (game continues, score is not persisted).
\end{itemize}

% ---- Module 4 ----
\subsection{Module 4: I/O Management}
I/O management defines \textit{how this computer talks to the world} --- screen, keyboard, and speaker.

\begin{table}[H]
\centering
\caption{I/O Device Mapping}
\begin{tabular}{lll}
\toprule
\textbf{Device} & \textbf{Terminal}         & \textbf{WebSocket}     \\
\midrule
Keyboard        & \texttt{getchar()} raw    & \texttt{recv()} + WS decode \\
Screen          & ANSI escape sequences     & WebSocket text frames  \\
Speaker         & WAV + platform player     & Web Audio API          \\
\bottomrule
\end{tabular}
\end{table}

\textbf{Audio --- cross-platform design:} \texttt{sound.c} generates square-wave WAV files at startup using pure integer arithmetic (no \texttt{<math.h>}). Playback uses a compile-time platform detection macro: \texttt{afplay} on macOS and \texttt{aplay} (ALSA) on Linux, with \texttt{paplay} (PulseAudio) as a secondary Linux fallback. A runtime availability check using \texttt{access()} confirms the binary exists before any \texttt{fork()} is attempted --- if no player is found the game runs silently with no crash or zombie processes.

\textbf{WebSocket I/O:} \texttt{keyboard.c} owns the receive path (\texttt{keyboard\_recv\_ws} decodes RFC\,6455 masked frames). \texttt{screen.c} owns the send path (\texttt{screen\_send\_ws} encodes WebSocket text frames). Game logic in \texttt{ws\_tetris.c} calls only these abstractions --- zero POSIX socket headers appear outside the HAL.

% ---- Module 5 ----
\subsection{Module 5: Error Handling \& Security}
Error handling ensures \textit{graceful behaviour when things go wrong}.

\begin{table}[H]
\centering
\caption{Error Handling Matrix}
\begin{tabular}{ll}
\toprule
\textbf{Error Category}      & \textbf{Handler}                        \\
\midrule
Division by zero             & \texttt{t\_div}/\texttt{t\_mod} return 0 safely \\
Memory OOM                   & \texttt{t\_alloc} returns NULL; caller checks \\
Board out-of-bounds          & \texttt{t\_in\_bounds()} guards all writes \\
Buffer overflow              & \texttt{t\_strncpy()} enforces max-length \\
Double-free                  & Block-list walk detects and ignores      \\
Invalid pointer free         & Range check + ownership walk rejects it  \\
SIGINT / SIGTERM             & Flag set; loop exits; terminal restored  \\
SIGWINCH                     & Flag set; full redraw on next frame      \\
Client disconnect            & Server re-accepts without restart        \\
Malformed JSON / WebSocket   & Frame silently dropped; loop continues   \\
Handshake failure            & Retry with next client connection        \\
Audio player missing         & Runtime check; game runs silently        \\
\bottomrule
\end{tabular}
\end{table}

\textbf{Secure halt:} \texttt{screen\_panic(msg)} is called on unrecoverable errors (e.g.\ allocation failure at boot). It restores the terminal cursor and raw-mode state, prints a \texttt{[KERNEL PANIC]} message to \texttt{stderr}, and calls \texttt{exit(1)} --- the shell is always left in a usable state.

% ---- Module 6 ----
\subsection{Module 6: Networking}
Networking defines \textit{how this computer talks to another computer}.

\textbf{Protocol:} WebSocket (RFC\,6455) over TCP on port 8080.

\textbf{Key features:}
\begin{itemize}
  \item \textbf{Persistent server socket:} survives page refresh and client reconnection. The outer accept loop re-accepts a new client without restarting the process.
  \item \textbf{Auto-reconnect client:} the React frontend retries with 1.5\,s exponential backoff on disconnect.
  \item \textbf{Full RFC\,6455 compliance:} SHA-1 handshake (implemented without \texttt{<openssl>}), masked frame decoding, FIN bit handling, close frame detection.
  \item \textbf{JSON state protocol:} full game state serialised every dirty tick (board, current/ghost/next/held pieces, score, leaderboard). Serialisation uses only \texttt{t\_itoa} and manual string building --- no \texttt{sprintf} or \texttt{<string.h>}.
  \item \textbf{Hardware abstraction:} all socket and POSIX network headers (\texttt{sys/socket.h}, \texttt{arpa/inet.h}) are confined exclusively to \texttt{screen.c} and \texttt{keyboard.c} --- zero networking code appears in game logic.
\end{itemize}

% ---- Module 7 ----
\subsection{Module 7: User Interface}
The user interface defines \textit{how a user interacts with the game}.

\textbf{Terminal UI:} ANSI-colored 10$\times$20 board with piece-type colors, ghost piece (shown as \texttt{-{}-}), next piece preview, hold piece display, HUD panel with score/level/lines/combo, game-over overlay showing the top-5 leaderboard, player name entry via \texttt{readLine()}, 4 block themes switchable with \texttt{T}, and background music.

\textbf{Web UI:} React\,18 + Vite with CRT scanline aesthetic, 4 selectable color themes (Neon, Retro, Mono, Pastel), hold piece, ghost piece, 7-bag randomizer, full SRS wall kicks (5-test per rotation), lock delay, DAS/ARR, combo system with live counter, level-up flash animation, top-10 leaderboard tab, mobile touch controls, SFX and music toggles, and auto-reconnect status indicator.

% ============================================================================
\section{Implementation Details}

\subsection{Custom Library Pipeline}

The six libraries form an explicit data-flow pipeline used in every game action:

\begin{lstlisting}[caption={Data Flow Pipeline}]
Input  -> keyboard.c (keyPressed / keyboard_recv_ws)
Parse  -> string.c   (t_strcmp, t_split, t_strlen)
State  -> memory.c   (t_alloc / t_dealloc)
Logic  -> math.c     (t_mul, t_mod, t_div, t_in_bounds)
Output -> screen.c   (ANSI escapes / screen_send_ws)
Audio  -> sound.c    (WAV gen / fork+exec player)
\end{lstlisting}

\textbf{Concrete example --- key dispatch:} \texttt{keyPressed()} returns a \texttt{char}; \texttt{key\_matches()} wraps it in a 2-byte string and calls \texttt{t\_strcmp()} for comparison; the matched action calls \texttt{t\_in\_bounds()} before any board write; the updated score is converted with \texttt{t\_itoa()} and sent via \texttt{screen\_render\_string()}. Every layer of the pipeline is exercised in a single keypress.

\subsection{Gameplay Mechanics}

\begin{table}[H]
\centering
\caption{Gameplay Feature Matrix}
\begin{tabular}{lcc}
\toprule
\textbf{Feature}        & \textbf{Terminal} & \textbf{Web}   \\
\midrule
7 Tetrominoes           & \checkmark        & \checkmark     \\
Ghost Piece             & \checkmark        & \checkmark     \\
Soft / Hard Drop        & \checkmark        & \checkmark     \\
Level Progression       & \checkmark        & \checkmark     \\
Sound Effects           & \checkmark        & \checkmark     \\
Background Music        & \checkmark        & \checkmark     \\
Player Name Entry       & \checkmark        & \checkmark     \\
Score Persistence       & \checkmark        & \checkmark     \\
Pause / Resume          & \checkmark        & \checkmark     \\
Hold Piece              & \checkmark        & \checkmark     \\
7-Bag Randomizer        & \checkmark        & \checkmark     \\
Lock Delay              & \checkmark        & \checkmark     \\
Combo System            & \checkmark        & \checkmark     \\
Leaderboard             & Top 5             & Top 10         \\
Theme Switcher          & 4 block styles    & 4 color themes \\
SRS Wall Kicks          & Basic (1-test)    & Full (5-test)  \\
DAS / ARR               & ---               & \checkmark     \\
Touch Controls          & ---               & \checkmark     \\
\bottomrule
\end{tabular}
\end{table}

\subsection{Technology Stack}

\begin{table}[H]
\centering
\caption{Technology Components}
\begin{tabular}{ll}
\toprule
\textbf{Component}  & \textbf{Technology}               \\
\midrule
Terminal backend    & C (gcc -Wall -Wextra -O2)         \\
WebSocket backend   & C (gcc -Wall -Wextra -O2)         \\
Frontend            & React 18 + Vite                   \\
Transport           & WebSocket RFC 6455 / TCP          \\
Audio (macOS)       & WAV + afplay                      \\
Audio (Linux)       & WAV + aplay / paplay              \\
Audio (Web)         & Web Audio API oscillators         \\
Persistence         & Plain text (pipe-delimited)       \\
Build               & GNU Make / npm                    \\
\bottomrule
\end{tabular}
\end{table}

% ============================================================================
\section{Results and Testing}

\subsection{Compilation}
Both build targets compile with zero warnings and zero errors under strict flags:
\begin{itemize}
  \item Terminal: \texttt{gcc -Wall -Wextra -O2 src/*.c -o tetris\_os}
  \item Backend:  \texttt{gcc -Wall -Wextra -O2 backend/ws\_tetris.c src/\{string,math,memory,screen,keyboard\}.c -o tetris\_ws}
\end{itemize}

\subsection{Demo}
A screen recording demonstrating both the terminal and web modes is available in the project repository at \url{https://github.com/Mehak261124/Tetris}. A screen recording demonstrating both modes is available in the \texttt{/demo/} folder of the repository.

\subsection{Edge Cases Verified}

\begin{table}[H]
\centering
\caption{Test Results}
\begin{tabular}{lcc}
\toprule
\textbf{Test Case}              & \textbf{Expected}   & \textbf{Result} \\
\midrule
Division by zero                & Returns 0           & Pass \\
\texttt{t\_alloc(0)}            & Returns NULL        & Pass \\
\texttt{t\_alloc} OOM           & Returns NULL        & Pass \\
\texttt{t\_dealloc(NULL)}       & No-op               & Pass \\
Double-free                     & Silently ignored    & Pass \\
Invalid pointer free            & Rejected            & Pass \\
SIGINT during gameplay          & Terminal restored   & Pass \\
SIGTERM during gameplay         & Terminal restored   & Pass \\
Terminal resize (SIGWINCH)      & Full redraw         & Pass \\
Terminal too small              & Warning displayed   & Pass \\
Page refresh (WS disconnect)    & Re-accepts          & Pass \\
Empty leaderboard file          & Starts fresh        & Pass \\
Malformed leaderboard line      & Line skipped        & Pass \\
Name overflow ($>$20 chars)     & Truncated safely    & Pass \\
Empty name input                & Defaults to Player  & Pass \\
Wall kick rotation (5 tests)    & All 5 attempted     & Pass \\
Hold twice per turn             & Second hold blocked & Pass \\
Malformed WebSocket frame       & Silently dropped    & Pass \\
WS handshake failure            & Retry next client   & Pass \\
Audio player not installed      & Silent, no crash    & Pass \\
\bottomrule
\end{tabular}
\end{table}

\subsection{Performance}
\begin{itemize}
  \item Game loop tick rate: 50\,ms (20\,Hz physics), 16\,ms render budget (~60\,fps)
  \item WebSocket frame size: 2--4\,KB per state update
  \item Memory footprint: 1\,MB arena allocated; game state occupies $<$2\,KB
  \item Frontend bundle: $<$100\,KB (production build)
\end{itemize}

% ============================================================================
\section{Known Issues}

Honest self-evaluation of current limitations:

\begin{itemize}
  \item \textbf{T-spin detection not implemented.} The \texttt{last\_was\_rotate} flag is tracked in \texttt{ws\_tetris.c} but the 4-corner detection logic and T-spin bonus scoring are not yet applied.
  \item \textbf{Terminal leaderboard limited to 5 entries.} The web backend stores 10 entries; the terminal mode stores 5. Both read the same file so the terminal silently discards entries 6--10 on load.
  \item \textbf{Background music does not loop seamlessly in the web client.} The melody restarts with a brief gap between repeats due to the \texttt{setTimeout}-based sequencer.
  \item \textbf{No DAS/ARR in terminal mode.} Rapid key repeat relies entirely on the OS keyboard repeat rate; the 170\,ms / 50\,ms DAS/ARR system is web-only.
  \item \textbf{Audio requires external binary on Linux.} \texttt{aplay} (ALSA) or \texttt{paplay} (PulseAudio) must be installed separately; the game degrades to silence if neither is present.
\end{itemize}

% ============================================================================
\section{Future Scope}

\begin{enumerate}
  \item \textbf{Multiplayer mode:} extend the WebSocket server to handle concurrent clients with garbage-line mechanics and a shared leaderboard.
  \item \textbf{T-spin detection:} implement 4-corner T-spin scoring using the existing \texttt{last\_was\_rotate} flag infrastructure.
  \item \textbf{Sprint mode:} timed 40-line clear challenge demonstrating advanced process scheduling (deadline-driven tick rate).
  \item \textbf{Replay system:} record timestamped inputs to file (file system extension) and replay them deterministically via the seeded LCG.
  \item \textbf{WebRTC signaling:} peer-to-peer multiplayer over the internet without a relay server.
  \item \textbf{Persistent sessions:} token-based session tracking so a web player's game survives a server restart.
  \item \textbf{DAS/ARR in terminal mode:} port the 170\,ms / 50\,ms repeat system from the web client to the keyboard.c raw-mode input path.
\end{enumerate}

% ============================================================================
\section{Conclusion}

The Tetris OS Simulator successfully demonstrates all seven fundamental OS module categories through a practical, interactive application. By replacing every standard library function with a custom implementation, the project provides a concrete, hands-on understanding of how operating systems manage processes, memory, I/O devices, file systems, and error conditions at the lowest level. The dual-mode architecture (terminal + WebSocket) showcases hardware abstraction principles, while the unified leaderboard file, real-time JSON protocol, and automatic reconnection demonstrate file system and networking implementations working together. Both build targets compile with zero warnings under \texttt{-Wall -Wextra}, comprehensive edge-case testing confirms robust error handling, and the full commit history at \url{https://github.com/Mehak261124/Tetris} documents the incremental development of each module.

% ============================================================================
\begin{thebibliography}{10}

\bibitem{c99}
ANSI/ISO 9899:1999, \textit{Programming Languages --- C}, International Organization for Standardization, 1999.

\bibitem{rfc6455}
I. Fette and A. Melnikov, ``The WebSocket Protocol,'' IETF RFC 6455, Dec. 2011. [Online]. Available: \url{https://tools.ietf.org/html/rfc6455}

\bibitem{sha1}
D. Eastlake and P. Jones, ``US Secure Hash Algorithm 1 (SHA-1),'' IETF RFC 3174, Sep. 2001.

\bibitem{guideline}
The Tetris Company, \textit{Tetris Guideline}, 2009. [Online]. Available: \url{https://tetris.wiki/Tetris_Guideline}

\bibitem{tanenbaum}
A. S. Tanenbaum and H. Bos, \textit{Modern Operating Systems}, 4th ed. Upper Saddle River, NJ: Pearson, 2014.

\bibitem{stallings}
W. Stallings, \textit{Operating Systems: Internals and Design Principles}, 9th ed. Upper Saddle River, NJ: Pearson, 2017.

\bibitem{kr}
B. W. Kernighan and D. M. Ritchie, \textit{The C Programming Language}, 2nd ed. Englewood Cliffs, NJ: Prentice Hall, 1988.

\bibitem{react}
React Team, ``React Documentation,'' Meta Open Source, 2024. [Online]. Available: \url{https://react.dev}

\bibitem{posix}
IEEE Std 1003.1-2017, \textit{The Open Group Base Specifications Issue 7}, IEEE, 2017.

\bibitem{knuth}
D. E. Knuth, \textit{The Art of Computer Programming, Vol. 2: Seminumerical Algorithms}, 3rd ed. Boston, MA: Addison-Wesley, 1997.

\end{thebibliography}

\end{document}